\chapter{Conclusiones y continuidad}\label{cap:conclusiones}

SavePoint entrega una plataforma funcional y un banco de pruebas de recomendación con un énfasis explícito en procedencia, trazabilidad y reproducibilidad. La contribución principal no es una única lista de juegos, sino el mecanismo que permite relacionar un ranking con sus señales, datos, parámetros, candidatas, métricas y artefactos. El autor diseñó y aplicó una metodología de ingeniería asistida por agentes, gobernada mediante GSD, organizada conceptualmente con Obsidian y validada mediante pruebas, revisiones y evidencias versionadas.

La publicación v15 muestra diferencias numéricas entre baselines, contenido, colaborativo e híbrido bajo sus condiciones fijadas. La evidencia no permite extrapolar esas cifras a usuarios reales ni convertir un único run de 79 usuarios sintéticos evaluables en superioridad científica general. Esa limitación es parte del resultado: delimita con precisión qué se ha demostrado y qué debe comprobarse después.

Las líneas de continuación son repetir el protocolo con nuevas semillas y cohortes, incorporar evaluación con usuarios y consentimiento, estudiar sensibilidad de pesos e hiperparámetros, completar las capturas y diagramas verificables, y consolidar el registro de esfuerzo y pruebas para la defensa. Cualquier extensión deberá crear un protocolo y artefacto nuevos, sin alterar retrospectivamente la evidencia v15.
